%% Согласно ГОСТ Р 7.0.11-2011:
%% 5.3.3 В заключении диссертации излагают итоги выполненного исследования, рекомендации, перспективы дальнейшей разработки темы.
%% 9.2.3 В заключении автореферата диссертации излагают итоги данного исследования, рекомендации и перспективы дальнейшей разработки темы.
\begin{enumerate}
  \item Разработана комплексная методика распознавания многокомпонентных
        структур по~их разрозненным частям, реализующая единый конвейер
        обработки из~четырех последовательных этапов: (1)~согласованное
        детектирование связанных компонентов структуры; (2)~раздельное
        извлечение дискриминативных эмбеддингов компонентов кодировщиками
        Vision Transformer с~функцией потерь ArcFace; (3)~вычисление
        гауссовых правдоподобий для каждого класса; (4)~принятие решения
        по~многоклассовому MAP-критерию. Методика обеспечивает поэтапное
        повышение качества на~каждой стадии при минимизации каскадного
        эффекта накопления ошибок и~устойчивое распознавание при неполноте
        данных; апробирована на~наиболее представительном частном случае~--
        реидентификации людей в~информационно-поисковых системах по~паре
        связанных компонентов~<<силуэт-лицо>> ($K=2$).

  \item Разработана байесовская модель объединения признаковых описаний
        компонентов многокомпонентной структуры и~решающее правило
        по~многоклассовому MAP-критерию с~автоматическим взвешиванием
        компонентов через ковариационные матрицы класса $\Sigma_y$
        в~гауссовых правдоподобиях. Модель обеспечивает адаптивный вклад
        каждого компонента без ручной подстройки весов и~корректное
        поведение решающего правила при произвольном множестве доступных
        компонентов (временном отсутствии или зашумлении части из~них).
        На~наиболее представительном частном случае реидентификации
        личности (набор CUHK03) достигнуты макроусредненные значения:
        Precision~-- 95,65\,\%, Recall~-- 95,54\,\%, $F_1$-мера~-- 94,31\,\%,
        Accuracy~-- 93,42\,\%; по~метрике mAP получены значения 95,65\,\% /
        98,3\,\% / 89,2\,\% для наборов CUHK03 / Market-1501 / MARS
        соответственно, что превосходит базовые способы линейного
        объединения признаков.

  \item Разработана модификация критерия обучения детектора (на~базе
        архитектуры YOLOv8) для согласованной локализации связанных
        компонентов структуры (внешнего и~вложенного), основанная
        на~введении дополнительного компонента функции потерь
        $\mathcal{L}_{\text{inner}}$ и~управляющего коэффициента
        $\lambda_{\text{inner}}$, совместно штрафующих ошибки локализации
        вложенного компонента в~границах внешнего. Экспериментально
        показано, что значения $\lambda_{\text{inner}}$ в~диапазоне
        2,0--2,7 обеспечивают наилучшие показатели по~детекции вложенного
        компонента (в~случае $K=2$~-- лица в~границах силуэта) без
        статистически значимой деградации качества детектирования базового
        внешнего класса.

  \item Подтверждена практическая значимость предложенной методики.
        Результаты работы внедрены и~используются в~деятельности
        ФГБУ~<<НИЦ~<<Институт имени Н.\,Е.\,Жуковского>>>>
        (сокращение времени принятия решения до~5\,\% в~системах
        визуального наблюдения).
\end{enumerate}
